\begin{abstract}
Recovery interventions on frozen robot policies are usually judged by the failures they repair, but the same intervention can also turn an episode the policy would have completed into a failure, and a net success gain hides that effect. We study this for a diffusion policy stacking cubes in simulation. When the policy is about to open its gripper over the target, an observation-only controller (Vision-CENTER V2) estimates the horizontal offset between the two cubes from a base and a wrist RGB camera and a short visual memory, and either abstains or makes a bounded horizontal move; the policy's action queue is then cleared and the policy re-plans. Each scene is branched from one simulator snapshot, so intervened and uninterrupted runs are paired, and rescues (failure to success) and harms (success to failure) are counted separately. After development, two pre-registered confirmatory stages used new scenes. On 400 calibration scenes per policy, harm was 1 of 348 and 0 of 337 baseline successes; exact binomial tests with Holm correction rejected a conditional harm rate of 5\% or more for both policies (family-wise error $\delta=0.05$). On 400 held-out test scenes per policy, run once, success rose from 350 to 379 and from 347 to 373 (29 rescues vs 0 harms, $p=1.9\times10^{-9}$; 27 vs 1, $p=1.1\times10^{-7}$). Much of the harm margin comes from scenes in which the controller abstains. We also report a failed contact-measurement check. The certificate is statistical and task-level, not a physical safety guarantee.
\end{abstract}

\begin{keyword}
recovery intervention \sep frozen visuomotor policy \sep paired evaluation \sep outcome harm \sep finite-sample certificate \sep diffusion policy \sep robotic manipulation
\end{keyword}
